% ------------------------------------------------------------
% Unified Altitude Model: A6 → A7
% Author: Borealis S. Hedling
% Date: October 2026
% Purpose: Formal LaTeX consolidation of dual/non-dual algebra,
%          A6/A7 glossaries, altitude boundaries, and A7 operators.
% ------------------------------------------------------------

\documentclass[11pt]{article}
\usepackage{amsmath, amssymb, amsthm}

\title{Unified Altitude Model: A6 Dual Algebra and A7 Non-Dual Algebra}
\author{Borealis S. Hedling}
\date{October 2026}

\begin{document}
\maketitle

% ------------------------------------------------------------
\section*{Terminology Notice}

The terms \emph{dual} and \emph{non-dual} in this document refer exclusively to
algebraic altitude regimes in the A6/A7 model. They have no relationship to
“dual-use” in regulatory, military, weapons, aerospace, or export-control
contexts.

Structural non-duality applies to both A6 and A7.  
Formal non-dual algebra applies only to A7.

% ------------------------------------------------------------
\section{Dual vs Non-Dual Algebra: Formal Definitions}

\subsection{Structural Non-Duality (A6 + A7)}
A structural design principle:
\begin{itemize}
  \item no primal/dual split,
  \item no adjoint geometry,
  \item no mirrored operator spaces,
  \item no kernel/surface duality,
  \item no collapse/continuity duality.
\end{itemize}

\subsection{Formal Dual Algebra (A6)}
A6 operators satisfy:
\begin{itemize}
  \item continuity,
  \item monotonicity,
  \item adjacency preservation,
  \item spectral boundedness,
  \item holonomy boundedness.
\end{itemize}

\subsection{Formal Non-Dual Algebra (A7)}
A7 operators satisfy:
\begin{itemize}
  \item discontinuity,
  \item non-monotonicity,
  \item adjacency fracture,
  \item spectral tearing,
  \item holonomy accumulation.
\end{itemize}

% ------------------------------------------------------------
\section{A6 Glossary: Dual Algebra}

\subsection{Spectral Continuity}
\[
|\delta(O(p))| \le \epsilon_{\mathrm{spec}}.
\]

\subsection{Modal Adjacency}
Adjacency graph $A(p,q)$ must remain preserved.

\subsection{Holonomy Drift}
\[
|\tau(p)| \le \epsilon_{\mathrm{holo}}.
\]

\subsection{Collapse Surface}
\[
p \in S \implies p \notin \Sigma.
\]

\subsection{Stability Envelope}
\[
E(p) \le \epsilon_S.
\]

\subsection{Unified Kernel}
\[
p \in K \implies O(p) \in K.
\]

% ------------------------------------------------------------
\section{A7 Glossary: Non-Dual Algebra}

\subsection{Divergence Invariant}
\[
I_{\mathrm{div}}(p) = E(p) - \epsilon_S.
\]

\subsection{Divergence Membrane}
\[
M = \partial S.
\]

\subsection{Divergence Kernel}
\[
K_{\mathrm{div}} = \{ p \mid I_{\mathrm{div}}(p) = 0 \}.
\]

\subsection{Divergence Stability Envelope}
\[
E_{\mathrm{div}}(p) = \max(0, I_{\mathrm{div}}(p)).
\]

\subsection{Divergence Collapse Surfaces}
\[
\Sigma_{\mathrm{div}} = \{ p \mid I_{\mathrm{div}}(p) = c \}.
\]

% ------------------------------------------------------------
\section{A6 $\leftrightarrow$ A7 Altitude Bridge}

\subsection{Altitude Boundary}
\[
A6: E(p) \le \epsilon_S,
\qquad
A7: p \in \partial S.
\]

\subsection{Operator Separation}
\[
O_{\mathrm{A6}} : S \to S,
\qquad
D_{\mathrm{A7}} : M \to M.
\]

\subsection{Kernel/Membrane Interaction}
\[
K_{\mathrm{div}} = \partial S.
\]

\subsection{Collapse Surface Interaction}
\[
T_p M + T_p \Sigma_{\mathrm{div}} = T_p X.
\]

% ------------------------------------------------------------
\section{Developer Note: Structural vs Formal Non-Duality}

A6 is structurally non-dual but formally dual.  
A7 is structurally non-dual and formally non-dual.

A7 extends A6 into the divergent regime without introducing dual structures.

% ------------------------------------------------------------
\section{A7 Non-Dual Projection Algebra}

\subsection{Overview}
Non-dual projection operators reconcile collapse-surface geometry at A7 altitude.

\subsection{Dual Projection Operators (A6)}
\[
P^2 = P.
\]
Preserve:
\begin{itemize}
  \item continuity,
  \item adjacency,
  \item spectral stability,
  \item holonomy bounds.
\end{itemize}

\subsection{Non-Dual Projection Operators (A7)}
\[
\Pi_{\mathrm{ND}} : M \to \Sigma.
\]

\[
\Pi_{\mathrm{ND}}(p) = p \quad \text{if } p \in \Sigma.
\]

Non-idempotent:
\[
\Pi_{\mathrm{ND}}^2 \ne \Pi_{\mathrm{ND}}.
\]

Non-linear:
\[
\Pi_{\mathrm{ND}}(ap + bq) \ne a\Pi_{\mathrm{ND}}(p) + b\Pi_{\mathrm{ND}}(q).
\]

Spectral discontinuity:
\[
\mathrm{Spec}(\Pi_{\mathrm{ND}}(p)) \not\approx \mathrm{Spec}(p).
\]

Holonomy breakage:
\[
H(\Pi_{\mathrm{ND}}(\gamma)) \not< \Theta_{\max}.
\]

\subsection{Domains and Codomains}
\[
\mathrm{Dom}(\Pi_{\mathrm{ND}}) = M,
\qquad
\mathrm{Cod}(\Pi_{\mathrm{ND}}) = \Sigma.
\]

\subsection{Weak Invariants}
Preserves:
\begin{itemize}
  \item collapse-surface membership,
  \item reconciliation locality,
  \item spectral tear closure,
  \item modal adjacency restoration.
\end{itemize}

Breaks:
\begin{itemize}
  \item continuity,
  \item determinism,
  \item kernel structure,
  \item spectral stability,
  \item holonomy bounds.
\end{itemize}

% ------------------------------------------------------------
\section{Conclusion}

This unified document formalizes the A6 dual algebra, the A7 non-dual algebra,
their altitude boundary, and the A7 non-dual projection operators. It provides a
complete mathematical foundation for Lean formalization and future structural AI
interpretability work.

% ------------------------------------------------------------
\section*{Provenance}

Generated by Borealis S. Hedling using the A6/A7 altitude model. All content is
simulation-safe, non-regulatory, and non-dual-use in the weapons or
export-control sense. This artifact is part of a reproducible research lineage.

\end{document}
